\begin{abstract}
Tomato maturity detectors are usually compared using one value: mean average
precision (mAP) on one dataset. A counting system deployed on an edge device is
also affected by factors that mAP does not record. These include the annotation
team, the visible fraction of each fruit, the confidence threshold used to
produce a count, the lighting, the model size, and the inference hardware. This
work measures the effect of each factor on YOLO11 detectors trained on two
public greenhouse datasets, Laboro Tomato and tomatOD. We map both datasets onto
the same three maturity stages and deploy the detectors with Core ML on the
Neural Engine of an iPhone~17~Pro. The datasets place the boundary between
half-ripe and fully ripe fruit at different colours: 98\% of tomatOD's fully
ripe boxes fall within the colour range that Laboro labels half-ripe. A detector
trained on both datasets retains its accuracy on each dataset and improves
transfer in both directions. As a larger fraction of each fruit is hidden,
detection recall falls from 0.884 to 0.057 at 70\% hidden. Among the fruit still
detected, the assigned stage is correct for 86\% to 94\% of them. The matching
rule and the appearance of the occluder each change recall by about ten points.
Selecting the confidence threshold reduces cross-dataset counting error from
3.29 to 1.34 fruit per image and stage without changing mAP. Across the five
YOLO11 sizes trained under one protocol, mAP stops improving at the medium
model, although package size and latency continue to increase. On the
smartphone, YOLO11x is the largest size that still runs at 30 frames per second,
with almost no margin. On a development laptop, the smallest model runs 3.85
times faster than on the smartphone, while the largest runs only 1.21 times
faster. Laptop timings therefore overstate the speed advantage of small models.
Three stops of over-exposure reduce the best model's recall by 0.205. This loss
is more than three times the 0.058 recall difference between the smallest and
largest models, and wider training-time colour augmentation does not recover
it. Reproducing a reported counting error therefore requires the annotation
source, occlusion and matching rules, confidence threshold, lighting conditions,
and deployment configuration to be reported.
\end{abstract}

\begin{IEEEkeywords}
agricultural computer vision, object detection, tomato ripeness, edge
inference, on-device deep learning, occlusion, domain shift, calibration,
illumination robustness, model scaling, Core ML
\end{IEEEkeywords}
